% GENERATED FILE. Edit canonical lessons, then run npm run generate:books.
% canonical-source-hash: fnv1a64:41bbe9ef6370ad89
% canonical-lessons: SA-C351-satkarma, SA-C351-vivekah, SA-C351-anuvartate, SA-C351-anusarati, SA-C351-maryada

\chapter{Right and Wrong}
\label{ch:sa-right-and-wrong}
\hlchaptermodality{351}

\begin{chapteropening}
\textbf{By the end of this chapter:} \emph{I can say what a good person should and should not do.}

Complete the last lesson of chapter 351: I can say what a good person should and should not do.
\end{chapteropening}

\section[\texorpdfstring{\sk{सत्कर्म} (satkarma)}{सत्कर्म (satkarma)}]{\sk{सत्कर्म} (satkarma) — a virtuous act}

\label{lesson:SA-C351-satkarma}



\begin{quote}
Before the new one: say the Sanskrit for gives back, then the Sanskrit for atonement.
\end{quote}

\subsection*{You'll want to know: \sk{सत्कर्म}}
\textbf{\sk{सत्कर्म}} — \emph{satkarma} — \textquotedblleft{}a virtuous act\textquotedblright{}.

\begin{quote}
\sk{सेवा} \sk{सत्कर्म} \sk{अस्ति।}
\end{quote}

\begin{quote}
\emph{sevā satkarma asti.}
\end{quote}

\begin{quote}
Service is a virtuous act.
\end{quote}

A neuter noun.

\begin{cousinweb}
\textbf{\sk{सत्कर्म}} (\emph{satkarma}) is built from \emph{sat}, \textquotedblleft{}good\textquotedblright{}, and \emph{karman}, \textquotedblleft{}an act\textquotedblright{}, from the root √kṛ; MW gives \emph{sat-karman} \textquotedblleft{}a good work, a virtuous act\textquotedblright{} (MW).
\end{cousinweb}

\begin{grammarlens}[title={What you've built}]
One more word for giving reasons and opinions.
\end{grammarlens}

\subsection*{Your turn}
\begin{itemize}
\raggedright
  \item \emph{Say it:} \emph{satkarma}
  \item \emph{Say it:} \emph{satkarma}, once more
  \item \emph{Say it:} \emph{sevā satkarma asti.} — Service is a virtuous act.
  \item \emph{Recall:} say the Sanskrit for gives back, then the Sanskrit for atonement, then say \emph{satkarma} again
\end{itemize}

\begin{tcolorbox}[breakable,colback=teal!4,colframe=teal!35!black,title={Before you move on}]
What does \sk{सत्कर्म} mean? (\textquotedblleft{}A virtuous act\textquotedblright{}.) Say it once more.
\end{tcolorbox}
\section[\texorpdfstring{\sk{विवेकः} (vivekaḥ)}{विवेकः (vivekaḥ)}]{\sk{विवेकः} (vivekaḥ) — discernment; good judgement}

\label{lesson:SA-C351-vivekah}



\begin{quote}
Before the new one: say the Sanskrit for atonement, then the Sanskrit for a virtuous act.
\end{quote}

\subsection*{You'll want to know: \sk{विवेकः}}
\textbf{\sk{विवेकः}} — \emph{vivekaḥ} — \textquotedblleft{}discernment; good judgement\textquotedblright{}.

\begin{quote}
\sk{विवेकः} \sk{सत्यं} \sk{मिथ्या} \sk{च} \sk{दर्शयति।}
\end{quote}

\begin{quote}
\emph{vivekaḥ satyaṃ mithyā ca darśayati.}
\end{quote}

\begin{quote}
Discernment shows what is true and what is false.
\end{quote}

A masculine noun.

\begin{cousinweb}
\textbf{\sk{विवेकः}} (\emph{vivekaḥ}) is built on \emph{vi} and the root √vic, \textquotedblleft{}to sift, to separate\textquotedblright{}; MW gives \emph{vi-veka} \textquotedblleft{}discrimination, distinction\textquotedblright{} and \textquotedblleft{}right judgement\textquotedblright{} (MW).
\end{cousinweb}

\begin{grammarlens}[title={What you've built}]
One more word for giving reasons and opinions.
\end{grammarlens}

\subsection*{Your turn}
\begin{itemize}
\raggedright
  \item \emph{Say it:} \emph{vivekaḥ}
  \item \emph{Say it:} \emph{vivekaḥ}, once more
  \item \emph{Say it:} \emph{vivekaḥ satyaṃ mithyā ca darśayati.} — Discernment shows what is true and what is false.
  \item \emph{Recall:} say the Sanskrit for atonement, then the Sanskrit for a virtuous act, then say \emph{vivekaḥ} again
\end{itemize}

\begin{tcolorbox}[breakable,colback=teal!4,colframe=teal!35!black,title={Before you move on}]
What does \sk{विवेकः} mean? (\textquotedblleft{}Discernment; good judgement\textquotedblright{}.) Say it once more.
\end{tcolorbox}
\section[\texorpdfstring{\sk{अनुवर्तते} (anuvartate)}{अनुवर्तते (anuvartate)}]{\sk{अनुवर्तते} (anuvartate) — conforms to; follows (a custom)}

\label{lesson:SA-C351-anuvartate}



\begin{quote}
Before the new one: say the Sanskrit for a virtuous act, then the Sanskrit for discernment.
\end{quote}

\subsection*{You'll want to know: \sk{अनुवर्तते}}
\textbf{\sk{अनुवर्तते}} — \emph{anuvartate} — \textquotedblleft{}conforms to; follows (a custom)\textquotedblright{}.

\begin{quote}
\sk{पुत्रः} \sk{धर्मस्य} \sk{मार्गम्} \sk{अनुवर्तते।}
\end{quote}

\begin{quote}
\emph{putraḥ dharmasya mārgam anuvartate.}
\end{quote}

\begin{quote}
The son follows the path of dharma.
\end{quote}

Past (he …): \sk{सः} \sk{धर्मम्} \sk{अनुवृत्तवान्}, \emph{saḥ dharmam anuvṛttavān}.

\begin{cousinweb}
\textbf{\sk{अनुवर्तते}} (\emph{anuvartate}) is built on \emph{anu}, \textquotedblleft{}after\textquotedblright{}, and the root √vṛt, \textquotedblleft{}to turn, to be\textquotedblright{}; MW gives \emph{anu-√vṛt} \textquotedblleft{}to follow\textquotedblright{} and \textquotedblleft{}to obey, to respect, to imitate\textquotedblright{} (MW).
\end{cousinweb}

\begin{grammarlens}[title={What you've built}]
One more word for giving reasons and opinions.
\end{grammarlens}

\subsection*{Your turn}
\begin{itemize}
\raggedright
  \item \emph{Say it:} \emph{anuvartate}
  \item \emph{Say it:} \emph{anuvartate}, once more
  \item \emph{Say it:} \emph{putraḥ dharmasya mārgam anuvartate.} — The son follows the path of dharma.
  \item \emph{Recall:} say the Sanskrit for a virtuous act, then the Sanskrit for discernment, then say \emph{anuvartate} again
\end{itemize}

\begin{tcolorbox}[breakable,colback=teal!4,colframe=teal!35!black,title={Before you move on}]
What does \sk{अनुवर्तते} mean? (\textquotedblleft{}Conforms to; follows (a custom)\textquotedblright{}.) Say it once more.
\end{tcolorbox}
\section[\texorpdfstring{\sk{अनुसरति} (anusarati)}{अनुसरति (anusarati)}]{\sk{अनुसरति} (anusarati) — abides by; keeps to}

\label{lesson:SA-C351-anusarati}



\begin{quote}
Before the new one: say the Sanskrit for discernment, then the Sanskrit for conforms to.
\end{quote}

\subsection*{You'll want to know: \sk{अनुसरति}}
\textbf{\sk{अनुसरति}} — \emph{anusarati} — \textquotedblleft{}abides by; keeps to\textquotedblright{}.

\begin{quote}
\sk{छात्रः} \sk{नियमान्} \sk{अनुसरति।}
\end{quote}

\begin{quote}
\emph{chātraḥ niyamān anusarati.}
\end{quote}

\begin{quote}
The student abides by the rules.
\end{quote}

Past (he …): \sk{सः} \sk{उपदेशम्} \sk{अनुसृतवान्}, \emph{saḥ upadeśam anusṛtavān}.

\begin{cousinweb}
\textbf{\sk{अनुसरति}} (\emph{anusarati}) is built on \emph{anu}, \textquotedblleft{}after\textquotedblright{}, and the root √sṛ, \textquotedblleft{}to run, to flow\textquotedblright{}; MW gives \emph{anu-sṛta} \textquotedblleft{}followed, conformed to\textquotedblright{} (MW).
\end{cousinweb}

\begin{grammarlens}[title={What you've built}]
One more word for giving reasons and opinions.
\end{grammarlens}

\subsection*{Your turn}
\begin{itemize}
\raggedright
  \item \emph{Say it:} \emph{anusarati}
  \item \emph{Say it:} \emph{anusarati}, once more
  \item \emph{Say it:} \emph{chātraḥ niyamān anusarati.} — The student abides by the rules.
  \item \emph{Recall:} say the Sanskrit for discernment, then the Sanskrit for conforms to, then say \emph{anusarati} again
\end{itemize}

\begin{tcolorbox}[breakable,colback=teal!4,colframe=teal!35!black,title={Before you move on}]
What does \sk{अनुसरति} mean? (\textquotedblleft{}Abides by; keeps to\textquotedblright{}.) Say it once more.
\end{tcolorbox}
\section[\texorpdfstring{\sk{मर्यादा} (maryādā)}{मर्यादा (maryādā)}]{\sk{मर्यादा} (maryādā) — the bounds of propriety}

\label{lesson:SA-C351-maryada}



\begin{quote}
Before the new one: say the Sanskrit for conforms to, then the Sanskrit for abides by.
\end{quote}

\subsection*{You'll want to know: \sk{मर्यादा}}
\textbf{\sk{मर्यादा}} — \emph{maryādā} — \textquotedblleft{}the bounds of propriety\textquotedblright{}.

\begin{quote}
\sk{बालः} \sk{मर्यादां} \sk{न} \sk{त्यजेत्।}
\end{quote}

\begin{quote}
\emph{bālaḥ maryādāṃ na tyajet.}
\end{quote}

\begin{quote}
A child should not overstep what is proper.
\end{quote}

A feminine noun.

\begin{cousinweb}
\textbf{\sk{मर्यादा}} (\emph{maryādā}) MW gives \emph{maryādā} \textquotedblleft{}a limit, a boundary\textquotedblright{} and \textquotedblleft{}the bounds of morality and propriety\textquotedblright{}; MW doubts whether it comes from \emph{maryā} with \emph{dā} or with \emph{ādā}, so its origin is uncertain (MW).
\end{cousinweb}

\begin{grammarlens}[title={What you've built}]
One more word for giving reasons and opinions.
\end{grammarlens}

\subsection*{Your turn}
\begin{itemize}
\raggedright
  \item \emph{Say it:} \emph{maryādā}
  \item \emph{Say it:} \emph{maryādā}, once more
  \item \emph{Say it:} \emph{bālaḥ maryādāṃ na tyajet.} — A child should not overstep what is proper.
  \item \emph{Recall:} say the Sanskrit for conforms to, then the Sanskrit for abides by, then say \emph{maryādā} again
\end{itemize}

\begin{tcolorbox}[breakable,colback=teal!4,colframe=teal!35!black,title={Before you move on}]
What does \sk{मर्यादा} mean? (\textquotedblleft{}The bounds of propriety\textquotedblright{}.) Say it once more.
\end{tcolorbox}
